\textbf{Phần 1. Trắc nghiệm lựa chọn}
\setcounterpageref{ex}{0}

\begin{ex} %[4H2N09]
Góc có số đo bằng $90^\circ$ được gọi là
\choice
{Góc nhọn}
{Góc tù}
{Góc bẹt}
{\True Góc vuông}
\loigiai{
Góc $90^\circ$ là góc vuông.\\
\textbf{Đáp án của câu hỏi là D}
}
\end{ex}

\begin{ex} %[4H2N09]
Góc có số đo lớn hơn $90^\circ$ và nhỏ hơn $180^\circ$ được gọi là
\choice
{Góc nhọn}
{\True Góc tù}
{Góc vuông}
{Góc bẹt}
\loigiai{
Đó là góc tù.\\
\textbf{Đáp án của câu hỏi là B}
}
\end{ex}

\begin{ex} %[4H2H09]
Lúc 9 giờ đúng, góc tạo bởi kim giờ và kim phút của chiếc đồng hồ là
\choice
{Góc nhọn}
{Góc tù}
{\True Góc vuông}
{Góc bẹt}
\loigiai{
Lúc 9 giờ đúng tạo thành góc vuông $90^\circ$.\\
\textbf{Đáp án của câu hỏi là C}
}
\end{ex}

\begin{ex} %[4H2H09]
Số đo của góc bẹt bằng
\choice
{\True $180^\circ$}
{$90^\circ$}
{$270^\circ$}
{$360^\circ$}
\loigiai{
Góc bẹt $= 180^\circ$.\\
\textbf{Đáp án của câu hỏi là A}
}
\end{ex}

\begin{ex} %[4H2H09]
Trong các hình sau, hình nào có 4 góc vuông?
\choice
{Hình bình hành}
{\True Hình chữ nhật}
{Hình thoi}
{Hình thang}
\loigiai{
Hình chữ nhật (và hình vuông) có 4 góc vuông.\\
\textbf{Đáp án của câu hỏi là B}
}
\end{ex}

\begin{ex} %[4H2H09]
Góc $75^\circ$ là
\choice
{Góc vuông}
{Góc tù}
{\True Góc nhọn}
{Góc bẹt}
\loigiai{
$75^\circ < 90^\circ$ là góc nhọn.\\
\textbf{Đáp án của câu hỏi là C}
}
\end{ex}

\begin{ex} %[4H2V09]
Lúc 8 giờ đúng, góc nhỏ tạo bởi kim giờ và kim phút có số đo là
\choice
{\True $120^\circ$}
{$90^\circ$}
{$150^\circ$}
{$180^\circ$}
\loigiai{
Khoảng cách giữa số 8 và số 12 là 4 khoảng giờ $\Rightarrow 4 \times 30^\circ = 120^\circ$.\\
\textbf{Đáp án của câu hỏi là A}
}
\end{ex}

\begin{ex} %[4H2V09]
Một góc tù có số đo lớn hơn góc nhọn $50^\circ$. Biết góc nhọn có số đo $65^\circ$. Số đo của góc tù đó là
\choice
{$105^\circ$}
{\True $115^\circ$}
{$125^\circ$}
{$135^\circ$}
\loigiai{
$65 + 50 = 115^\circ$.\\
\textbf{Đáp án của câu hỏi là B}
}
\end{ex}

\begin{ex} %[4H2V09]
Khi kim phút quay từ số 12 đến số 6 thì kim phút đã quét một góc có số đo là
\choice
{$90^\circ$}
{$120^\circ$}
{$360^\circ$}
{\True $180^\circ$}
\loigiai{
Từ số 12 đến số 6 là nửa vòng tròn, tương ứng góc bẹt $180^\circ$.\\
\textbf{Đáp án của câu hỏi là D}
}
\end{ex}

\begin{ex} %[4H2V09]
Góc tạo bởi hai cạnh kề nhau của một chiếc ê-ke tam giác vuông cân gồm có
\choice
{\True 1 góc vuông và 2 góc nhọn $45^\circ$}
{3 góc vuông}
{2 góc vuông và 1 góc nhọn}
{3 góc nhọn}
\loigiai{
Ê-ke tam giác vuông cân có 1 góc vuông $90^\circ$ và 2 góc nhọn $45^\circ$.\\
\textbf{Đáp án của câu hỏi là A}
}
\end{ex}

\begin{ex} %[4H2V09]
Cho hai góc kề bù nhau (tổng bằng $180^\circ$), trong đó một góc có số đo gấp đôi góc kia. Số đo góc tù là
\choice
{$100^\circ$}
{$135^\circ$}
{\True $120^\circ$}
{$150^\circ$}
\loigiai{
Tổng số phần: $1 + 2 = 3$ phần.\\
Góc tù: $180 : 3 \times 2 = 120^\circ$.\\
\textbf{Đáp án của câu hỏi là C}
}
\end{ex}

\begin{ex} %[4H2C09]
Một chiếc quạt giấy khi xòe hết cỡ tạo thành một hình bán nguyệt (nửa hình tròn). Góc tạo bởi hai nan ngoài cùng của quạt là
\choice
{Góc vuông}
{Góc tù}
{Góc đầy ($360^\circ$)}
{\True Góc bẹt ($180^\circ$)}
\loigiai{
Hình nửa đường tròn ứng với góc bẹt $180^\circ$.\\
\textbf{Đáp án của câu hỏi là D}
}
\end{ex}

\textbf{Phần 2. Trắc nghiệm đúng sai}
\setcounterpageref{ex}{0}

\begin{ex} %[4H2H09]
Xét tính đúng sai của các nhận định về góc:
\choiceTF
{\True Dùng thước đo góc để xác định số đo của góc.}
{\True Góc nhọn luôn có số đo bé hơn $90^\circ$.}
{Hai góc nhọn cộng lại luôn tạo thành một góc tù.}
{\True Góc bẹt có số đo bằng tổng số đo của hai góc vuông.}
\loigiai{
\begin{itemize}
    \item Dùng thước đo góc: Đúng.
    \item Góc nhọn $< 90^\circ$: Đúng.
    \item Hai góc nhọn có thể cộng lại thành góc nhọn (ví dụ $20^\circ + 30^\circ = 50^\circ$). Sai.
    \item $90^\circ + 90^\circ = 180^\circ$ (góc bẹt). Đúng.
\end{itemize}
\begin{itemize}
    \item a) \textbf{Đúng}.
    \item b) \textbf{Đúng}.
    \item c) \textbf{Sai}.
    \item d) \textbf{Đúng}.
\end{itemize}
}
\end{ex}

\begin{ex} %[4H2V09]
Một hình tứ giác $ABCD$ có góc $A = 90^\circ$, góc $B = 90^\circ$, góc $C = 120^\circ$. Xét tính đúng sai:
\choiceTF
{\True Góc đỉnh $A$ và đỉnh $B$ là hai góc vuông.}
{\True Góc đỉnh $C$ là góc tù.}
{Góc đỉnh $D$ có số đo là $80^\circ$.}
{\True Góc đỉnh $D$ có số đo là $60^\circ$ (là góc nhọn).}
\loigiai{
\begin{itemize}
    \item Góc $A, B$ vuông: Đúng.
    \item Góc $C = 120^\circ$ tù: Đúng.
    \item Góc $D = 360^\circ - (90^\circ + 90^\circ + 120^\circ) = 60^\circ \neq 80^\circ$. Ý 3 Sai, Ý 4 Đúng.
\end{itemize}
\begin{itemize}
    \item a) \textbf{Đúng}.
    \item b) \textbf{Đúng}.
    \item c) \textbf{Sai}.
    \item d) \textbf{Đúng}.
\end{itemize}
}
\end{ex}

\textbf{Phần 3. Trả lời ngắn}
\setcounterpageref{ex}{0}

\begin{ex} %[4H2H09]
Tính số đo góc tạo bởi hai kim đồng hồ lúc 10 giờ đúng (theo độ, góc nhỏ hơn $180^\circ$).
\shortans{60}
\loigiai{
\textbf{Đáp số:} $60$\\
Khoảng cách là 2 giờ $\Rightarrow 2 \times 30^\circ = 60^\circ$.
}
\end{ex}

\begin{ex} %[4H2H09]
Góc vuông gấp mấy lần góc $30^\circ$?
\shortans{3}
\loigiai{
\textbf{Đáp số:} $3$\\
$90 : 30 = 3$ lần.
}
\end{ex}

\begin{ex} %[4H2V09]
Một góc tù có số đo bằng $135^\circ$. Hỏi góc này lớn hơn góc vuông bao nhiêu độ?
\shortans{45}
\loigiai{
\textbf{Đáp số:} $45$\\
$135 - 90 = 45^\circ$.
}
\end{ex}

\begin{ex} %[4H2V09]
Một góc bẹt được chia thành 3 góc bằng nhau. Tính số đo của mỗi góc đó (theo độ).
\shortans{60}
\loigiai{
\textbf{Đáp số:} $60$\\
$180 : 3 = 60^\circ$.
}
\end{ex}

\begin{ex} %[4H2V09]
Một hình tam giác có 1 góc vuông và 1 góc nhọn bằng $40^\circ$. Tính số đo góc nhọn còn lại (theo độ).
\shortans{50}
\loigiai{
\textbf{Đáp số:} $50$\\
$90 - 40 = 50^\circ$.
}
\end{ex}

\begin{ex} %[4H2C09]
Lúc mấy giờ đúng trong khoảng từ 6 giờ đến 12 giờ thì hai kim đồng hồ tạo thành một góc vuông? (Chỉ ghi số giờ).
\shortans{9}
\loigiai{
\textbf{Đáp số:} $9$\\
Lúc 9 giờ đúng hai kim tạo thành góc vuông.
}
\end{ex}
